\section{Product-belief-first：把行为信念写进界面、数据与后训练}

第二条路径从一句更强的产品主张开始：系统应该怎样与人共同工作？这类目标往往无法从单一 benchmark 推导。团队需要先构造交互原型，让用户暴露“哪里像合作者、哪里像自动补全”，再把这些差异转成训练样本、reward 与 eval。产品 belief 不是凌驾于证据之上的审美，而是等待被操作化和反驳的 hypothesis。

\lecturefigure{slide-23-product-belief-vision.jpg}{路径二：从 product belief 与 vision 出发，再让模型做到}{00:13:39--00:14:28}

\begin{importantbox}{Product belief 的可验收写法}
“模型要更聪明”“体验要更自然”都不可直接验收。可操作的 belief 应包含：目标用户与任务、模型应采取/避免的行为、用户保留的控制权、失败时的恢复路径、成功 evidence 和成本边界。例如“成为写作 collaborator”至少意味着模型能读取选区、提出局部修改、说明改动、保持作者声音，并允许逐步撤销。
\end{importantbox}

\subsection{Language as a high-end product}

讲者把 language 本身视为高端产品材料。语言不是界面上的填充物：语气、信息密度、结构、直接性、引用和不确定性表达会改变用户如何理解世界。新闻产品的例子说明，同一 context block 可以在 story page、首页 carousel 和 bottom sheet 中呈现不同形态；底层事实共享，表面组织必须适配场景。

\lecturefigure{slide-24-language-high-end-product.jpg}{Language as a high-end product：信息组织本身就是产品行为}{00:14:28--00:14:36}

\lecturefigure{slide-25-context-block-multiple-surfaces.jpg}{共享 context block：同一事实跨 surface 保持一致}{00:14:36--00:14:45}

\lecturefigure{slide-26-adaptive-presentation-surfaces.jpg}{同一信息在内联 carousel、首页与 bottom sheet 中自适应呈现}{00:14:45--00:15:39}

可以把共享语义表示记为 $z$，用户情境为 $u$，surface 为 $s$，最终呈现为
\[
y_s=g_s(z,u).
\]
$g_s$ 负责篇幅、布局和交互方式，但不得悄悄改变 $z$ 中的事实。工程上应把 factual core、presentation policy 和 user state 分层；否则一个 surface 的编辑可能在另一个 surface 产生不一致，或模型为了适应短界面删掉关键限定。

\begin{knowledgebox}{首次使用：Single source of truth}
Single source of truth 不是要求所有界面显示同一句话，而是让多个 surface 引用同一可追踪事实对象。模型可以摘要、重排和改写，但每个结论应能回到共享来源；用户在一个位置更新事实后，其他位置应明确同步或标注版本。
\end{knowledgebox}

\subsection{从 Writing IDE 到 micro-personalization}

Writing IDE 例子把语言模型放进作者工作流，而不是要求作者把整篇文稿复制进聊天框。真正难点不是生成一句漂亮文字，而是理解选区、上下文、作者意图和改动范围。Micro-personalization 则进一步要求系统记住“这个用户怎样工作”，但个性化不能演变成不可见的操控。

\lecturefigure{slide-27-writing-ide-gpt3.jpg}{Writing IDE：模型在文档上下文中提供局部协作}{00:15:39--00:15:50}

\lecturefigure{slide-28-language-compressed-knowledge.jpg}{语言是人类知识的高度压缩 artifact，也是界面创新的材料}{00:15:50--00:16:17}

\lecturefigure{slide-29-claude-micro-personalization.jpg}{Claude.ai micro-personalization：根据历史主题组织后续入口}{00:16:17--00:16:20}

若个性化状态为 $m_u$，一次更新可写为
\[
m_u' = U(m_u,e_t,\gamma),
\]
其中 $e_t$ 是本次交互 evidence，$\gamma$ 是保留强度。产品必须回答：哪些事件可写入、保存多久、用户如何查看/删除、错误记忆如何回滚。若 $U$ 完全隐形，模型可能把一次偶然偏好固化成长期画像。

\begin{warningbox}{个性化的三个误区}
第一，把高频点击当作稳定偏好；第二，在用户不知道时跨任务迁移敏感信息；第三，为了“贴心”而减少观点多样性。Micro-personalization 应同时测 relevance、surprise、user control 与 privacy，而不是只测点击率。
\end{warningbox}

\teachervoice{00:13:39--00:16:18，Karina 从新闻产品、Writing IDE 和语言的压缩性说明 product belief 如何先于模型行为。老师强调：真正的 interface innovation 往往要求训练模型理解新的交互对象，而不是只在现有模型外面画 UI。}

\subsection{Virtual teammate：异步、多人与工具访问}

“AI teammate”常被滥用为营销词。讲者给出的 Claude in Slack 原型有更具体的行为边界：异步对话、多人协作、自动任务、定期总结频道和工具访问。它仍不是完整同事，但至少把 teammate 从人格形容词改写成可观察的工作协议。

\lecturefigure{slide-30-claude-slack-product-spec.jpg}{Claude in Slack：从频道历史生成 product spec、目标与 OKR}{00:16:20--00:17:50}

\begin{knowledgebox}{术语消化：Assistant、Tool、Agent、Teammate}
\begin{tabularx}{\textwidth}{L{0.17\textwidth} X X}
\toprule
角色 & 核心机制 & 主要验收问题 \\
\midrule
Assistant & 对请求生成响应 & 单轮帮助是否正确、清楚 \\
Tool & 执行窄操作 & 输入输出是否确定、可审计 \\
Agent & 维护状态并调用工具 & 权限、计划、恢复、停止条件 \\
Teammate & 在多人流程中持续协调 & 责任边界、共享上下文、交接与社会判断 \\
\bottomrule
\end{tabularx}
同一个模型可以承担不同角色；角色由环境与权限定义，不由聊天语气定义。
\end{knowledgebox}

一个多人频道中的 teammate 至少要维护共享状态 $s_t$、个人可见状态 $p_i$ 与权限集合 $\mathcal{A}_i$。可执行动作必须满足
\[
a_t\in \mathcal{A}(s_t,p_i,\text{policy}),
\]
而不是模型“认为有帮助”就执行。总结频道、写 spec 和修改 issue 的风险等级不同，应有不同确认门槛。

\subsection{训练模型成为 collaborator}

Canvas 不是把文本编辑器接到模型上就完成。讲者回顾的训练工作包括目标数据、context 使用、文档格式、editing suggestion、rewriting、formatting 和 output quality。图中的柱状结果意味着团队把“协作”拆成多个可重复任务，再迭代数据与后训练；它并不证明一个总分能代表所有写作场景。

\lecturefigure{slide-31-human-ai-flexible-affordances.jpg}{Human--AI collaborative affordance：界面应随任务扩展}{00:17:50--00:19:36}

\lecturefigure{slide-32-training-collaborator.jpg}{训练模型成为 collaborator：多类编辑行为与质量比较}{00:18:21--00:19:41}

\begin{importantbox}{Collaborator 的最小行为 contract}
\begin{enumerate}
  \item 读取当前 artifact、选区、作者目标与禁止改动范围；
  \item 先提出局部 plan，再生成可比较的 diff；
  \item 保持未授权区域和作者 voice；
  \item 对事实性改动给出处或不确定性；
  \item 允许接受、拒绝、局部修改与完整 rollback；
  \item 把用户反馈写成下一轮可复现 eval，而不是只做在线适配。
\end{enumerate}
\end{importantbox}

\begin{lstlisting}[caption={带控制权的文档协作循环}]
state = load(document, selection, user_goal, permissions)
proposal = model.plan_and_edit(state)
diff = isolate_changes(proposal, protected_regions=state.locked)
evidence = verify_facts_and_style(diff, state.sources)
decision = user.review(diff, evidence)
if decision.accepted:
    commit(diff)
else:
    record_failure(state, diff, decision.reason)
\end{lstlisting}

\teachervoice{00:17:50--00:20:25，课堂追问“模型成为 collaborator 到底是什么意思”，并指出 Canvas 需要专门训练 context use、formatting、rewriting 与 editing suggestion。实践经验：好的协作界面背后是行为数据与 eval，不是一个更大的输入框。}

\subsection{Modular tool composition：Tasks 不是单个 prompt}

ChatGPT Tasks 画面把定时触发、生成文稿、再次编辑和保存组合起来。这里的系统不只生成文本，还需要 scheduler、artifact store、权限与失败恢复。把多个模块连接起来会产生新的 state mismatch：任务可能在旧上下文上运行，生成结果可能覆盖用户修改，重复执行可能制造冲突。

\lecturefigure{slide-33-chatgpt-tasks-composition.jpg}{ChatGPT Tasks：定时任务、Canvas artifact 与后续编辑的模块组合}{00:19:41--00:20:25}

\begin{knowledgebox}{首次使用：Idempotency 与 rollback}
\term{Idempotency（幂等性）}指同一动作重复执行不会产生额外副作用；定时任务尤其需要 idempotency key，避免网络重试重复发文或重复修改文件。\term{Rollback（回滚）}则要求系统能恢复到已知版本。两者是 agent/tool 产品的基础，不应等到模型出错后再补。
\end{knowledgebox}

\subsection{本章小结}

Product-belief-first 的核心是把“希望模型怎样合作”拆成界面对象、权限、数据与 eval。语言产品、共享 context、个性化、Slack teammate、Canvas collaborator 和 Tasks 都说明：越接近真实工作，越需要版本、权限、diff、rollback 和用户控制。下一章进入本讲最扎实的案例——如何把模糊 model behavior 变成可信 eval。

